% !TEX root = ../thesis.tex

\chapter{Vyhodnotenie}\label{ch:evaluation}

Kapitola sumarizuje výsledky experimentov pre tri modelové rodiny (Phi-3 Mini, Gemma 3 4B, Qwen3 4B) v troch režimoch: \emph{base}, \emph{fine-tuned} (SFT) a \emph{DPO}. Hodnotenie bolo realizované jednotným evaluačným skriptom s rovnakými dekódovacími parametrami a rovnakou sadou doménových otázok, aby bolo možné korektne porovnať vplyv doménového doučenia na kvalitu a praktické parametre inferencie.

\section{Experimentálne nastavenie a metriky}

Použité metriky pokrývajú kvalitu výstupov aj prevádzkovú použiteľnosť:
\begin{itemize}
    \item \textbf{G-Eval} (LLM-as-a-judge): \emph{relevance, coherence, fluency, factual accuracy},
    \item \textbf{perplexita} na testovacích vzorkách (nižšia je lepšia),
    \item \textbf{TTFT} a \textbf{TPS} ako latencia a priepustnosť generovania,
    \item \textbf{VRAM} po načítaní modelu a pri kontexte 1024 tokenov.
\end{itemize}

\section{Kvalitatívne výsledky (G-Eval)}

Tabuľka~\ref{tab:geval} uvádza priemerné G-Eval skóre pre všetky varianty.

\begin{table}[h]
\centering
\small
\begin{tabular}{lcccc}
\hline
Model & Relevance & Coherence & Fluency & Factual accuracy \\
\hline
Phi-3 Base & 8.83 & 8.71 & 9.21 & 8.54 \\
Phi-3 SFT & 9.38 & 9.00 & 9.24 & 9.00 \\
Phi-3 DPO & 9.32 & 9.11 & 9.21 & 9.43 \\
Gemma-3 Base & 9.76 & 9.43 & 9.28 & 9.28 \\
Gemma-3 SFT & 9.74 & 9.37 & 9.41 & 9.48 \\
Gemma-3 DPO & 9.70 & 9.41 & 9.48 & 9.63 \\
Qwen3 Base & 9.83 & 9.38 & 9.34 & 9.24 \\
Qwen3 SFT & 9.72 & 9.38 & 9.31 & 9.53 \\
Qwen3 DPO & 9.93 & 9.43 & 9.29 & 9.33 \\
\hline
\end{tabular}
\caption{G-Eval skóre (priemery) pre jednotlivé modely.}
\label{tab:geval}
\end{table}

Z výsledkov vyplýva:
\begin{itemize}
    \item \textbf{Phi-3:} doménové SFT prináša konzistentné zlepšenie vo všetkých dimenziách; DPO ďalej zvyšuje najmä \emph{factual accuracy} (z 8.54 na 9.43).
    \item \textbf{Gemma-3:} základný model je už veľmi silný; SFT a DPO zvyšujú plynulosť a faktickú správnosť, pričom najvyššiu \emph{factual accuracy} dosahuje Gemma-3 DPO (9.63).
    \item \textbf{Qwen3:} DPO má najvyššiu \emph{relevance} (9.93) a mierne zlepšuje koherenciu, SFT zvyšuje \emph{factual accuracy} (9.53), ale mierne znižuje fluency.
\end{itemize}

\section{Perplexita a stručnosť odpovedí}

Perplexita (Tabuľka~\ref{tab:ppl}) po SFT klesá vo všetkých rodinách, čo indikuje lepšie prispôsobenie na doménový formát. DPO naopak perplexitu oproti \emph{base} takmer nemení. Výrazný je aj pokles priemernej dĺžky odpovede: všetky doučené varianty generujú podstatne kratšie odpovede, čo je v súlade s cieľom používať stručný, štruktúrovaný výstup.

\begin{table}[h]
\centering
\small
\begin{tabular}{lcc}
\hline
Model & Perplexita $\downarrow$ & Priemerná dĺžka odpovede (tokeny) \\
\hline
Phi-3 Base & 7.138 & 465.3 \\
Phi-3 SFT & 5.167 & 223.4 \\
Phi-3 DPO & 7.128 & 250.9 \\
Gemma-3 Base & 16.807 & 901.7 \\
Gemma-3 SFT & 8.193 & 232.0 \\
Gemma-3 DPO & 16.786 & 249.5 \\
Qwen3 Base & 14.423 & 658.6 \\
Qwen3 SFT & 13.027 & 265.9 \\
Qwen3 DPO & 14.361 & 272.2 \\
\hline
\end{tabular}
\caption{Perplexita a priemerná dĺžka odpovedí.}
\label{tab:ppl}
\end{table}

\section{Prevádzkové metriky a pamäťové nároky}

Prevádzkové metriky ukazujú typický kompromis medzi kvalitou a rýchlosťou. \emph{Base} varianty majú najnižší \texttt{TTFT}, zatiaľ čo SFT/DPO mierne zvyšujú latenciu prvej odpovede a znižujú \texttt{TPS}. Výhodou je kratšia priemerná odpoveď, čo v praxi znižuje celkový čas generovania.

\begin{table}[h]
\centering
\small
\begin{tabular}{lcccc}
\hline
Model & TTFT [s] $\downarrow$ & TPS $\uparrow$ & VRAM load [MB] & VRAM pri 1024 [MB] \\
\hline
Phi-3 Base & 0.0827 & 9.41 & 858.3 & 1236.8 \\
Phi-3 SFT & 0.1141 & 7.63 & 2224.4 & 2594.8 \\
Phi-3 DPO & 0.1139 & 7.46 & 3993.3 & 4364.5 \\
Gemma-3 Base & 0.0892 & 6.94 & 3314.4 & 3973.6 \\
Gemma-3 SFT & 0.1138 & 7.71 & 3319.2 & 3979.8 \\
Gemma-3 DPO & 0.1132 & 7.62 & 3318.0 & 3978.8 \\
Qwen3 Base & 0.0829 & 7.72 & 2734.4 & 3036.1 \\
Qwen3 SFT & 0.1166 & 7.18 & 2738.0 & 3039.8 \\
Qwen3 DPO & 0.1173 & 7.15 & 2735.4 & 3037.2 \\
\hline
\end{tabular}
\caption{Prevádzkové metriky inferencie a pamäťové nároky.}
\label{tab:perf}
\end{table}

Z pohľadu VRAM je výrazný rozdiel najmä pri Phi-3: DPO verzia vyžaduje približne 4~GB VRAM, zatiaľ čo \emph{base} len 0.86~GB. Gemma-3 a Qwen3 majú medzi režimami prakticky stabilnú pamäťovú stopu, čo zjednodušuje ich nasadenie.

\section{Zhrnutie a odporúčaný kompromis}

Doménové doučenie prináša najväčší kvalitatívny prínos pri Phi-3, kde SFT a DPO výrazne zvyšujú relevanciu a faktickú správnosť. Pri Gemma-3 a Qwen3 je prínos jemnejší, ale prejavuje sa najmä v \emph{factual accuracy} a v skracovaní odpovedí. Z hľadiska kompromisu kvalita vs. náročnosť nasadenia je najvyváženejšou voľbou \textbf{Qwen3 SFT}: zachováva veľmi vysoké G-Eval skóre, zvyšuje faktickú správnosť na 9.53 a zároveň patrí k pamäťovo najúspornejším 4B modelom. Ak je prioritou maximálna faktická správnosť, najlepší výsledok poskytuje \textbf{Gemma-3 DPO} (9.63), a ak je kľúčovým obmedzením VRAM a latencia, najľahšie nasaditeľný je \textbf{Phi-3 Base}.
